\documentclass[12pt, a4paper]{article}

% Codificação e idioma
\usepackage[utf8]{inputenc}
\usepackage[T1]{fontenc}
\usepackage[brazil]{babel}

% Geometria e margens
% Define as margens. Você pode ajustar os valores 'left', 'right', 'top' e 'bottom'.
% Valores menores (e.g., 1.5cm) dão mais espaço.
\usepackage[
    top=1.5cm,
    bottom=1.5cm,
    left=1.5cm,
    right=1.5cm,
    a4paper,
    % Para cortar o cabeçalho/rodapé
    includeheadfoot
]{geometry}

% Espaçamento de linhas
% Você pode usar 'singlespacing' para o menor espaçamento possível.
% Ou, se precisar de um espaçamento personalizado, use \setstretch{0.9} por exemplo.
\usepackage{setspace}
\singlespacing

% Seções e Títulos
% Reduz o espaço antes e depois dos títulos das seções.
\usepackage{titlesec}
\titlespacing{\section}{0pt}{1.5ex}{1ex}
\titlespacing{\subsection}{0pt}{1.2ex}{1ex}
\titlespacing{\subsubsection}{0pt}{1ex}{0.8ex}

% Tabelas
\usepackage{booktabs}

% Matemática e gráficos
\usepackage{amsmath}
\usepackage{graphicx}
\usepackage{float}

% Tabela de conteúdo e bibliografia
\usepackage[nottoc]{tocbibind}


\usepackage{parskip}
\setlength{\parindent}{0pt} % removendo paragrafo

\newcommand{\myparagraph}[1]{\paragraph{#1}\mbox{}} % adicionando uma quebra de linha

\begin{document}

\section{Rede Neural para Aproximação de Soluções em Problemas de Otimização Linear}

\section*{Problema}
Problemas de Programação Linear (PL) são onipresentes em logística, produção, telecom e finanças. Métodos exatos (e.g., Simplex) são a referência, porém têm custo de execução por instância. O trabalho propõe uma rede neural para aprender a aproximar a solução ótima de PLs pequenos.


\subsection{Fonte dos Dados}

Neste projeto, os dados são gerados artificialmente (sintéticos). Cada instância é um problema de programação linear com duas variáveis e duas restrições. As soluções ótimas são obtidas via biblioteca PuLP, que serve para compara com o treinamento da rede neural.

\subsection{Objetivo}
Treinar uma rede neural para receber como entrada os coeficientes do problema e prever como saída a solução ótima aproximada. Dessa forma, o modelo pode atuar como um "solucionador rápido", útil em aplicações de tempo real.

\subsection{Descrição do problma}

Dado $[c_1, c_2, a_1, a_2, a_3, a_4, b_1, b_2]$, preveja a solução aproximada $[x_1^*, x_2^*]$ de um PL do tipo:
\begin{align*}
\text{maximizar } & c_1x_1 + c_2x_2 \\
\text{sujeito a } & a_1x_1 + a_2x_2 \leq b_1 \\
& a_3x_1 + a_4x_2 \leq b_2 \\
& x_1, x_2 \geq 0
\end{align*}

\subsection{Resolução do problema}

\begin{tabular}{|p{0.25\textwidth}|p{0.75\textwidth}|}
\hline
\textbf{Etapa} & \textbf{Descrição} \\
\hline
\textbf{1. Geração de Dados} & Gerar $N$ problemas de Programação Linear (PL) aleatórios e resolvê-los com a biblioteca PuLP para obter a solução ótima. \\
\hline
\textbf{2. Preparação de Dados} & Preparar os dados (coeficientes e soluções) convertendo-os em tensores. \\
\hline
\textbf{3. Treinamento da Rede} & Treinar uma rede neural com a seguinte configuração: \\
& \begin{itemize}
    \item \textbf{Arquitetura:} Rede $8 \to 32 \to 16 \to 2$
    \item \textbf{Função de Ativação:} ReLU
    \item \textbf{Otimizador:} Adam
    \item \textbf{Função de Perda:} MSE
    \item \textbf{Framework:} PyTorch
\end{itemize} \\
\hline
\end{tabular}

\subsection{Limitações}
A rede é treinada com dados sintéticos e limitada a problemas de PL com duas variáveis de decisão e duas restrições, logo não pode generalizar problemas maiores.


\end{document}